% GENERATED FILE. Edit canonical lessons, then run npm run generate:books.
% canonical-source-hash: fnv1a64:138ca703300a98d5
% canonical-lessons: LA-C420-consto, LA-C420-auctio, LA-C420-liceor, LA-C420-emptor, LA-C420-venditor, LA-R420-review-413-420

\chapter{Buying and Selling}
\label{ch:la-buying-and-selling}
\hlchaptermodality{420}

\begin{chapteropening}
\textbf{By the end of this chapter:} \emph{I can say what something costs, talk about an auction and bidding, and name the buyer and the seller.}

Complete the last lesson of chapter 420: I can say what something costs, talk about an auction and bidding, and name the buyer and the seller.
\end{chapteropening}

\section[cōnstō, cōnstāre]{cōnstō, cōnstāre — to cost; to stand firm; to be agreed}

\label{lesson:LA-C420-consto}



\begin{quote}
Before the new one: say the Latin for a written bond, then the Latin for a moneylender.
\end{quote}

\subsection*{You'll want to know: cōnstō, cōnstāre}
\textbf{cōnstō, cōnstāre} — \textquotedblleft{}to cost; to stand firm; to be agreed\textquotedblright{}.

\begin{quote}
Haec vīlla magnō cōnstitit.
\end{quote}

\begin{quote}
\emph{This country house cost a great deal.}
\end{quote}

\emph{cōnstō, cōnstāre, cōnstitī}. The price is in the ablative: \emph{magnō}, at a high price, \emph{parvō}, cheaply. \emph{Cōnstat}, it is agreed. English \emph{cost} comes from it through French.

\begin{cousinweb}
\emph{cōnstō} is built from \emph{con-} and \emph{stō} \textquotedblleft{}to stand\textquotedblright{} (EtymDB; Lewis \& Short): to stand firm, then to stand at a price. English \emph{cost} comes from it through French.
\end{cousinweb}

\begin{grammarlens}[title={What you've built}]
One more word for telling, explaining and describing at B1.
\end{grammarlens}

\subsection*{Your turn}
\begin{itemize}
\raggedright
  \item \emph{Say it:} \emph{cōnstō, cōnstāre}
  \item \emph{Say it:} \emph{cōnstō, cōnstāre}, once more
  \item \emph{Say it:} a sentence of your own with \emph{cōnstō, cōnstāre}
  \item \emph{Recall:} say the Latin for a written bond, then the Latin for a moneylender, then say \emph{cōnstō, cōnstāre} again
\end{itemize}

\begin{tcolorbox}[breakable,colback=teal!4,colframe=teal!35!black,title={Before you move on}]
What does cōnstō, cōnstāre mean? (\textquotedblleft{}To cost; to stand firm; to be agreed\textquotedblright{}.) Say it once more.
\end{tcolorbox}
\section[auctiō, auctiōnis]{auctiō, auctiōnis — an auction}

\label{lesson:LA-C420-auctio}



\begin{quote}
Before the new one: say the Latin for a moneylender, then the Latin for to cost.
\end{quote}

\subsection*{You'll want to know: auctiō, auctiōnis}
\textbf{auctiō, auctiōnis} — \textquotedblleft{}an auction\textquotedblright{}.

\begin{quote}
Bona eius in auctiōne vendita sunt.
\end{quote}

\begin{quote}
\emph{His goods were sold at an auction.}
\end{quote}

\emph{auctiō, auctiōnis} f., built on \emph{augeō}, to increase: the price goes up and up. English \emph{auction}.

\begin{cousinweb}
\emph{auctiō} is the noun built on \emph{augeō} \textquotedblleft{}to increase\textquotedblright{} (EtymDB; Lewis \& Short): a sale where the price keeps rising. English \emph{auction} comes from it.
\end{cousinweb}

\begin{grammarlens}[title={What you've built}]
One more word for telling, explaining and describing at B1.
\end{grammarlens}

\subsection*{Your turn}
\begin{itemize}
\raggedright
  \item \emph{Say it:} \emph{auctiō, auctiōnis}
  \item \emph{Say it:} \emph{auctiō, auctiōnis}, once more
  \item \emph{Say it:} a sentence of your own with \emph{auctiō, auctiōnis}
  \item \emph{Recall:} say the Latin for a moneylender, then the Latin for to cost, then say \emph{auctiō, auctiōnis} again
\end{itemize}

\begin{tcolorbox}[breakable,colback=teal!4,colframe=teal!35!black,title={Before you move on}]
What does auctiō, auctiōnis mean? (\textquotedblleft{}An auction\textquotedblright{}.) Say it once more.
\end{tcolorbox}
\section[liceor, licērī]{liceor, licērī — to bid (at an auction)}

\label{lesson:LA-C420-liceor}



\begin{quote}
Before the new one: say the Latin for to cost, then the Latin for an auction.
\end{quote}

\subsection*{You'll want to know: liceor, licērī}
\textbf{liceor, licērī} — \textquotedblleft{}to bid (at an auction)\textquotedblright{}.

\begin{quote}
Hunc equum centum dēnāriīs licitus est.
\end{quote}

\begin{quote}
\emph{He bid a hundred denarii for this horse.}
\end{quote}

Deponent: \emph{liceor, licērī, licitus sum}. The thing bid for goes in the accusative and the price in the ablative. It belongs with \emph{licet}: to name what you are allowed to pay.

\begin{cousinweb}
\emph{liceor} belongs with \emph{liceō} \textquotedblleft{}to be for sale, to be valued\textquotedblright{} and \emph{licet} \textquotedblleft{}it is allowed\textquotedblright{} (EtymDB; Lewis \& Short): to put a value on a thing.
\end{cousinweb}

\begin{grammarlens}[title={What you've built}]
One more word for telling, explaining and describing at B1.
\end{grammarlens}

\subsection*{Your turn}
\begin{itemize}
\raggedright
  \item \emph{Say it:} \emph{liceor, licērī}
  \item \emph{Say it:} \emph{liceor, licērī}, once more
  \item \emph{Say it:} a sentence of your own with \emph{liceor, licērī}
  \item \emph{Recall:} say the Latin for to cost, then the Latin for an auction, then say \emph{liceor, licērī} again
\end{itemize}

\begin{tcolorbox}[breakable,colback=teal!4,colframe=teal!35!black,title={Before you move on}]
What does liceor, licērī mean? (\textquotedblleft{}To bid (at an auction)\textquotedblright{}.) Say it once more.
\end{tcolorbox}
\section[emptor, emptōris]{emptor, emptōris — a buyer}

\label{lesson:LA-C420-emptor}



\begin{quote}
Before the new one: say the Latin for an auction, then the Latin for to bid.
\end{quote}

\subsection*{You'll want to know: emptor, emptōris}
\textbf{emptor, emptōris} — \textquotedblleft{}a buyer\textquotedblright{}.

\begin{quote}
Emptor pretium nimis magnum esse dīxit.
\end{quote}

\begin{quote}
\emph{The buyer said the price was too high.}
\end{quote}

\emph{emptor, emptōris} m., built on \emph{emō}, to buy. \emph{Caveat emptor}: let the buyer beware.

\begin{cousinweb}
\emph{emptor} is built on \emph{emō} \textquotedblleft{}to buy\textquotedblright{} (Lewis \& Short), from its supine \emph{emptum}, with the ending \emph{-tor} of the one who does it.
\end{cousinweb}

\begin{grammarlens}[title={What you've built}]
One more word for telling, explaining and describing at B1.
\end{grammarlens}

\subsection*{Your turn}
\begin{itemize}
\raggedright
  \item \emph{Say it:} \emph{emptor, emptōris}
  \item \emph{Say it:} \emph{emptor, emptōris}, once more
  \item \emph{Say it:} a sentence of your own with \emph{emptor, emptōris}
  \item \emph{Recall:} say the Latin for an auction, then the Latin for to bid, then say \emph{emptor, emptōris} again
\end{itemize}

\begin{tcolorbox}[breakable,colback=teal!4,colframe=teal!35!black,title={Before you move on}]
What does emptor, emptōris mean? (\textquotedblleft{}A buyer\textquotedblright{}.) Say it once more.
\end{tcolorbox}
\section[venditor, venditōris]{venditor, venditōris — a seller}

\label{lesson:LA-C420-venditor}



\begin{quote}
Before the new one: say the Latin for to bid, then the Latin for a buyer.
\end{quote}

\subsection*{You'll want to know: venditor, venditōris}
\textbf{venditor, venditōris} — \textquotedblleft{}a seller\textquotedblright{}.

\begin{quote}
Venditor equum laudābat.
\end{quote}

\begin{quote}
\emph{The seller was praising the horse.}
\end{quote}

\emph{venditor, venditōris} m., built on \emph{vendō}, to sell. English \emph{vendor}.

\begin{cousinweb}
\emph{venditor} is built on \emph{vendō} \textquotedblleft{}to sell\textquotedblright{} (EtymDB; Lewis \& Short), from its supine \emph{venditum}. English \emph{vendor} comes from the same verb.
\end{cousinweb}

\begin{grammarlens}[title={What you've built}]
One more word for telling, explaining and describing at B1.
\end{grammarlens}

\subsection*{Your turn}
\begin{itemize}
\raggedright
  \item \emph{Say it:} \emph{venditor, venditōris}
  \item \emph{Say it:} \emph{venditor, venditōris}, once more
  \item \emph{Say it:} a sentence of your own with \emph{venditor, venditōris}
  \item \emph{Recall:} say the Latin for to bid, then the Latin for a buyer, then say \emph{venditor, venditōris} again
\end{itemize}

\begin{tcolorbox}[breakable,colback=teal!4,colframe=teal!35!black,title={Before you move on}]
What does venditor, venditōris mean? (\textquotedblleft{}A seller\textquotedblright{}.) Say it once more.
\end{tcolorbox}
\section[(dialogue)]{Words from chapters 413–420 — a review}

\label{lesson:LA-R420-review-413-420}



\begin{quote}
Say the Latin for to condemn and a conviction.
\end{quote}

\subsection*{Your turn}
Say these aloud:

\begin{itemize}
\raggedright
  \item to condemn, a conviction, an acquittal, to beg to be spared, defence in court
  \item a wrong, a disgraceful act, divine law, a crime against heaven, unfair
  \item to punish, to send away, to post for sale, a whip, custody
  \item justice, fairness, to obey, to restrain, to claim
  \item possession, one's own, belonging to another, ownership, formal purchase
  \item debt, a debtor, a creditor, interest, borrowed
  \item to pay in full, interest charged, a pledge, a written bond, a moneylender
  \item to cost, an auction, to bid, a buyer, a seller
\end{itemize}

\begin{tcolorbox}[breakable,colback=teal!4,colframe=teal!35!black,title={Before you move on}]
To condemn? (\textbf{condemnō, condemnāre}.) A conviction? (\textbf{damnātiō, damnātiōnis}.) An acquittal? (\textbf{absolūtiō, absolūtiōnis}.) To beg to be spared? (\textbf{dēprecor, dēprecārī}.) Defence in court? (\textbf{patrōcinium, patrōciniī}.) A wrong? (\textbf{iniūria, iniūriae}.) A disgraceful act? (\textbf{flāgitium, flāgitiī}.) Divine law? (\textbf{fās}.) A crime against heaven? (\textbf{nefās}.) Unfair? (\textbf{inīquus, inīqua, inīquum}.) To punish? (\textbf{castīgō, castīgāre}.) To send away? (\textbf{relēgō, relēgāre}.) To post for sale? (\textbf{prōscrībō, prōscrībere}.) A whip? (\textbf{flagellum, flagellī}.) Custody? (\textbf{custōdia, custōdiae}.) Justice? (\textbf{iūstitia, iūstitiae}.) Fairness? (\textbf{aequitās, aequitātis}.) To obey? (\textbf{obtemperō, obtemperāre}.) To restrain? (\textbf{coerceō, coercēre}.) To claim? (\textbf{vindicō, vindicāre}.) Possession? (\textbf{possessiō, possessiōnis}.) One's own? (\textbf{proprius, propria, proprium}.) Belonging to another? (\textbf{aliēnus, aliēna, aliēnum}.) Ownership? (\textbf{dominium, dominiī}.) Formal purchase? (\textbf{mancipium, mancipiī}.) Debt? (\textbf{aes aliēnum, aeris aliēnī}.) A debtor? (\textbf{dēbitor, dēbitōris}.) A creditor? (\textbf{crēditor, crēditōris}.) Interest? (\textbf{fēnus, fēnoris}.) Borrowed? (\textbf{mūtuus, mūtua, mūtuum}.) To pay in full? (\textbf{persolvō, persolvere}.) Interest charged? (\textbf{ūsūra, ūsūrae}.) A pledge? (\textbf{pignus, pignoris}.) A written bond? (\textbf{syngrapha, syngraphae}.) A moneylender? (\textbf{fēnerātor, fēnerātōris}.) To cost? (\textbf{cōnstō, cōnstāre}.) An auction? (\textbf{auctiō, auctiōnis}.) To bid? (\textbf{liceor, licērī}.) A buyer? (\textbf{emptor, emptōris}.) A seller? (\textbf{venditor, venditōris}.)
\end{tcolorbox}
